% 图 15.2 —— 手工重建（2026-10-06 第二版：**按原书实测重新锚定**）
%
% 上一版的问题（用户逐条点出）：
%   ① 开区间括号画在带**外面** ⇒ 括起来的坐标轴范围**超出**了 U 的实际区间，拓扑上讲不通
%   ② U 花括号被 `U` 标签**盖住** —— 根因是竖带下边我画到 595，原书只到 574（多伸 21px）
%   ③ 整体几何左偏：y 轴 93（应 106）、竖带 288/434（应 293/437）、横带 277/445（应 281/444）
%   ④ 阴影是 **96 条逐条摆的线**（描摹思维）—— 改回「一条带 = 一个 clip + 平行线族」
%
% ★★ 几何一律取自原书 `images/15.2.png`（**与本画布同框 925×641**，可直接量）：
%     y 轴 x=106      x 轴 y=531
%     横带 π₂⁻¹(V)：x 61..760, y 281..444
%     竖带 π₁⁻¹(U)：x 293..437, y 90..574
%     区间标签 U 在 x 轴上 (293,437)、V 在 y 轴上 (281,444)
%     花括号：V 在 x 37..57（尖朝左）；U 在 y 585..601（尖朝下）；`U` 标签 y 611..632
\documentclass[12pt]{standalone}
\usepackage{fontspec}
\setsansfont{TeX Gyre Heros}
\newfontfamily\mfallback{DejaVu Sans}
\usepackage{amsmath}
\usepackage{tikz}
\usetikzlibrary{arrows.meta,decorations.pathreplacing}
\begin{document}
\begin{tikzpicture}[x=1pt, y=-1pt, line width=7.0pt, line cap=round, line join=round]

  %% ── 实测几何 ──
  \def\Yax{106}   \def\Xax{531}                                  % 坐标轴
  \def\HBx{61}  \def\HBr{760} \def\HBy{281} \def\HBb{444}       % 横带 π₂⁻¹(V)
  \def\VBx{293} \def\VBr{437} \def\VBy{90}  \def\VBb{574}       % 竖带 π₁⁻¹(U)
  \def\d{19.1}                                                   % 阴影线沿 x 的间距

  %% ════ 坐标轴（一条命令）════
  \draw (\Yax,0) -- (\Yax,\Xax) -- (925,\Xax);

  %% ════ 两个带：**只画两条长边**，短端是开的（原书如此，不是矩形！）════
  \draw (\HBx,\HBy) -- (\HBr,\HBy);      % π₂⁻¹(V) 上边
  \draw (\HBx,\HBb) -- (\HBr,\HBb);      % π₂⁻¹(V) 下边
  \draw (\VBx,\VBy) -- (\VBx,\VBb);      % π₁⁻¹(U) 左边
  \draw (\VBr,\VBy) -- (\VBr,\VBb);      % π₁⁻¹(U) 右边

  %% ════ 阴影：45° 平行线族，每块用 clip 裁进带里（**一条带 = 一个元素**）════
  % ★ 2026-10-06 更正方向（用户指出的事实性错误，4× 放大原书核实）：
  %   **横带 = `\`（左上→右下）**、**竖带 = `/`（右上→左下）** —— 我原来两条带都画成 `/`，
  %   所以交叠区还得另补反向线；方向改对后，**交叠区的交叉网由两块阴影自然叠出**，不必特判。
  % ---- 横带 `\`（46 条）----
  \begin{scope}\clip (\HBx,\HBy) rectangle (\HBr,\HBb);
    \foreach \k in {-9,...,36}
      \draw[line width=4.5pt] ({\HBx+\k*\d},\HBy) -- ++(190,190);
  \end{scope}
  % ---- 竖带 `/`（34 条）★ 段长必须 ≳ 带的对角（484），190 够不到上半段 ----
  \begin{scope}\clip (\VBx,\VBy) rectangle (\VBr,\VBb);
    \foreach \k in {-25,...,7}
      \draw[line width=4.5pt] ({\VBx+\k*\d},\VBb) -- ++(600,-600);
  \end{scope}

  %% ════ 坐标轴上的开区间括号（**蓝色**；标 U、V 是开区间）════
  %  放在区间**里面**：弧的外缘正好落在区间端点上 ⇒ 括起来的范围 = U / V 本身
  %  ★ 尺寸按本项目**既有实测口径**（fig-drawing「brace-shape-not-tikz-decoration」记的 15.1 端帽）：
  %    圆心落在轴线上、半径 ≈26、扫 ≈130°（我初版 R=16、扫 80° ⇒ 用户："括号小了"）
  \draw[blue, line width=3.8pt] (319,\Xax) ++(115:26) arc[start angle=115, end angle=245, radius=26];   % x 轴左端 (
  \draw[blue, line width=3.8pt] (411,\Xax) ++(-65:26) arc[start angle=-65, end angle=65, radius=26];    % x 轴右端 )
  \draw[blue, line width=3.8pt] (\Yax,307) ++(-155:26) arc[start angle=-155, end angle=-25, radius=26]; % y 轴上端
  \draw[blue, line width=3.8pt] (\Yax,418) ++(25:26) arc[start angle=25, end angle=155, radius=26];     % y 轴下端

  %% ════ 花括号（各一条命令；尖朝**外**）════
  \draw[decorate,decoration={brace,amplitude=20pt},line width=4.0pt] (57,\HBb) -- (57,\HBy);   % 左：V（尖朝左，37..57）
  \draw[decorate,decoration={brace,amplitude=16pt},line width=4.0pt] (\VBr,585) -- (\VBx,585);  % 下：U（尖朝下，585..601）

  %% ════ 标签 ════
  % ★ 整条写成一个节点：上下标交给排版引擎（拆成 7 个字形节点会因"数字=帽高/π=x高/括号=括号高"口径不同而错位）
  \node[anchor=center,inner sep=0pt,fill=white] at (373.5,54)
    {\fontsize{30}{36}\selectfont
     $\text{\textsf{\textbf{\textit{π}}}}_{\text{\textsf{\textbf{1}}}}%
      ^{\text{\textsf{\textbf{\mfallback −}}}\text{\textsf{\textbf{1}}}}%
      \text{\textsf{\textbf{(\textit{U})}}}$};
  \node[anchor=center,inner sep=0pt,fill=white] at (822,360)
    {\fontsize{30}{36}\selectfont
     $\text{\textsf{\textbf{\textit{π}}}}_{\text{\textsf{\textbf{2}}}}%
      ^{\text{\textsf{\textbf{\mfallback −}}}\text{\textsf{\textbf{1}}}}%
      \text{\textsf{\textbf{(\textit{V})}}}$};
  \node[anchor=center,inner sep=0pt,fill=white,font=\fontsize{29.6}{37.0}\selectfont] at (22.8,363.0) {\textsf{\textbf{\textit{V}}}};
  \node[anchor=center,inner sep=0pt,fill=white,font=\fontsize{38.1}{47.6}\selectfont] at (815.8,563.5) {\textsf{\textbf{\textit{x}}}};
  \node[anchor=center,inner sep=0pt,fill=white,font=\fontsize{29.6}{37.0}\selectfont] at (359.0,621.0) {\textsf{\textbf{\textit{U}}}};

  % 钉死画布
  \pgfresetboundingbox
  \path[use as bounding box] (0,0) rectangle (925,641);
\end{tikzpicture}
\end{document}
